\section{HiGFlow: o solver}

Quarenta e dois fontes, 54\,457 linhas --- tr\^es vezes e meia a HiGTree. A maior
parte disso s\~ao \emph{modelos}, n\~ao infraestrutura.

\subsection{Um m\'odulo de passo por tipo de escoamento}

O \texttt{flowtype} da configura\c{c}\~ao seleciona qual passo roda:

\pgfplotstabletypeset[
  string type,
  columns/tipodeescoamento/.style={column name={\texttt{flowtype}}},
  columns/modulo/.style={column name={m\'odulo}},
  columns/linhas/.style={column name={linhas}},
  every head row/.style={before row=\toprule, after row=\midrule},
  every last row/.style={after row=\bottomrule},
]{\dat{higflow-passos.dat}}

S\~ao 19\,900 linhas s\'o em passos de escoamento. O maior deles ---
suspens\~oes, com 3\,413 --- \'e maior que \texttt{domain.c}, o n\'ucleo da
HiGTree.

\subsection{E os acoplamentos entre eles}

\pgfplotstabletypeset[
  string type,
  columns/modulo/.style={column name={m\'odulo}},
  columns/linhas/.style={column name={linhas}},
  columns/acopla/.style={column name={acopla}},
  every head row/.style={before row=\toprule, after row=\midrule},
  every last row/.style={after row=\bottomrule},
]{\dat{higflow-acoplados.dat}}

Os dois \'ultimos t\^em 31 e 28 linhas: s\~ao despachantes, n\~ao
implementa\c{c}\~oes. A combinat\'oria de f\'isicas est\'a declarada mais do que
realizada, e vale saber disso antes de contar com ela.

\subsection{Modelos constitutivos}

\pgfplotstabletypeset[
  string type,
  columns/familia/.style={column name={fam\'ilia}},
  columns/modelos/.style={column name={modelos}},
  every head row/.style={before row=\toprule, after row=\midrule},
  every last row/.style={after row=\bottomrule},
]{\dat{modelos.dat}}

\subsection{Discretiza\c{c}\~oes, por eixo}

\pgfplotstabletypeset[
  string type,
  columns/eixo/.style={column name={eixo}},
  columns/opcoes/.style={column name={op\c{c}\~oes}},
  every head row/.style={before row=\toprule, after row=\midrule},
  every last row/.style={after row=\bottomrule},
]{\dat{discretizacoes.dat}}

O produto cartesiano dessas op\c{c}\~oes \'e grande, e nada garante que todas as
combina\c{c}\~oes tenham sido exercitadas --- esta \'e uma lista do que o c\'odigo
\emph{aceita}, n\~ao do que foi verificado.

\subsection{Tratamento de interface}

\pgfplotstabletypeset[
  string type,
  columns/familia/.style={column name={fam\'ilia}},
  columns/modulos/.style={column name={m\'odulos}},
  columns/papel/.style={column name={papel}},
  every head row/.style={before row=\toprule, after row=\midrule},
  every last row/.style={after row=\bottomrule},
]{\dat{higflow-interface.dat}}

Onze m\'odulos de VOF, contra um de front-tracking e um de fronteira imersa. A
assimetria conta a hist\'oria do sistema: o VOF \'e o caminho maduro, com
variantes de reconstru\c{c}\~ao e de curvatura acumuladas ao longo de anos; o
front-tracking \'e recente e tem uma implementa\c{c}\~ao s\'o.

\subsection{Entrada e sa\'ida: o m\'odulo mais improv\'avel}

\texttt{hig-flow-io.c} tem \textbf{10\,972 linhas} --- \SI{20}{\percent} de todo o
HiGFlow, e mais que a soma de \texttt{domain.c} e \texttt{pdomain.c}. Ele l\^e as
configura\c{c}\~oes em YAML, grava e rel\^e o estado, e escreve VTK. Cada
combina\c{c}\~ao de \texttt{flowtype} com dimens\~ao tem o seu par de
leitura/escrita, e \'e da\'i que vem o tamanho.

\'E o m\'odulo que ningu\'em planeja e que cresce sozinho.
